\documentclass[10pt,a4paper]{article}
\usepackage{fontspec,amsmath,amssymb,graphicx}
\usepackage[margin=18mm]{geometry}
\setmainfont{Times New Roman}
\setlength{\parindent}{0pt}
\setlength{\parskip}{5pt}
\allowdisplaybreaks
\sloppy
\emergencystretch=2em
\begin{document}
\section*{BT\_3\_5\_3\_05 — Lời giải bổ sung}
\subsection*{Kiến thức và trình tự giải}
\textbf{Bước 1 — Chọn năng lượng theo toán tử.} Với sóng dùng $E=\frac12\int(u_t^2+a^2u_x^2)dx$. Với Klein–Gordon có thêm $u$, thêm $\frac12\int u^2dx$. Với dầm $u_{tt}=-a^2u_{xxxx}$ dùng $E=\frac12\int(u_t^2+a^2u_{xx}^2)dx$, không thay $u_{xx}$ bằng $u_x$.

\textbf{Bước 2 — Đạo hàm và tích phân từng phần.} Với $u_{tt}-a^2u_{xx}+\gamma u_t+\kappa u=F$, năng lượng có hạng $\kappa u^2$ cho


\[
E'=a^2[u_tu_x]_0^L-\gamma\int u_t^2dx+\int Fu_tdx.
\]


Với dầm, tích phân hai lần cho $E'=a^2[u_{tx}u_{xx}-u_tu_{xxx}]_0^L$. Viết riêng từng hạng biên rồi dùng điều kiện biên và đạo hàm theo thời gian của điều kiện đó để cho chúng bằng không.

\textbf{Bước 3 — Suy hệ quả.} Không cản, không nguồn thì $E$ bảo toàn. Có cản $\gamma>0$ thì $E'\le0$; nếu $E$ không có hệ số $1/2$, $E'=-2\gamma\int u_t^2\ge-2\gamma E$, suy $e^{-2\gamma t}E(0)\le E(t)\le E(0)$. Với hiệu hai nghiệm cùng dữ liệu, $E(0)=0$ nên có duy nhất.

\textbf{Bước 4 — Tách biến nếu phải tìm nghiệm cụ thể.} Sóng có cản $\gamma=1$ cho $T''+T'+a^2\lambda T=0$. Đặt $T=e^{-t/2}Y$, thu $Y''+(a^2\lambda-1/4)Y=0$, $Y(0)=A$, $Y'(0)=B+A/2$. Phân biệt dao động, nghiệm kép và hàm hyperbolic. Dầm cho $X''''=k^4X$, $T''+a^2k^4T=0$; dùng đủ bốn điều kiện biên để tìm phương trình của $kL$.

\textbf{Bước 5 — Kiểm tra hệ số và khẳng định của đề.} Năng lượng đầu chứa cả $g^2$ và các đạo hàm của $f$. Thử $f=0$, $g\ne0$ để phát hiện công thức bỏ vận tốc. Với dầm, nghiệm $k=0$ phải giải trực tiếp bằng đa thức bậc ba; nghiệm $0$ của phương trình lượng giác–hyperbolic có thể chỉ là nghiệm giả do cơ sở suy biến.

\textbf{Chi tiết phép triệt tiêu tích phân.} Đạo hàm động năng cho $\int u_tu_{tt}$; dùng PDE thay $u_{tt}$. Với sóng, $\int u_tu_{xx}=[u_tu_x]-\int u_{tx}u_x$, và tích phân cuối triệt tiêu đạo hàm thế năng $\int u_xu_{xt}$. Với dầm,


\[
\int u_tu_{xxxx}=[u_tu_{xxx}-u_{tx}u_{xx}]+\int u_{txx}u_{xx}.
\]


Tích phân cuối triệt tiêu đạo hàm $\frac12\int u_{xx}^2$ sau khi nhân hệ số và giữ đúng dấu PDE. Viết hạng biên trước, sau đó mới dùng điều kiện của từng đầu.

\textbf{Bước 1 — Tách biến và chọn bài toán trị riêng theo biên.} Đặt $u=X(x)T(t)$ trong $u_{tt}=a^2u_{xx}$, suy $X''+\lambda X=0$, $T''+a^2\lambda T=0$. Tích phân từng phần cho $\lambda\int X^2=\int (X')^2$ khi hạng biên bằng không; do đó $\lambda\ge0$. Giải cả trường hợp $\lambda=0$ trước khi xét $\lambda=k^2>0$.

\textbf{Bước 2 — Tìm hệ hàm riêng.} Trên $(0,L)$: Dirichlet hai đầu dùng $\sin(n\pi x/L)$, $n\ge1$; Neumann hai đầu dùng $1$ và $\cos(n\pi x/L)$; Dirichlet–Neumann dùng $\sin((n+1/2)\pi x/L)$, $n\ge0$; Neumann–Dirichlet dùng cos của các tần số nửa nguyên ấy. Chuẩn bình phương của mode khác không là $L/2$, của mode hằng là $L$.

\textbf{Bước 3 — Chiếu cả hai dữ liệu.} Với mode khác không đặt


\[
A_n=\frac2L\int_0^L fX_n\,dx,\quad B_n=\frac2L\int_0^L gX_n\,dx,\quad \omega_n=ak_n.
\]


Tích phân từng khoảng nếu dữ liệu chia đoạn; dùng tích phân từng phần hoặc công thức tích–tổng khi là đa thức/lượng giác. Hệ số $B_n$ là hệ số của vận tốc, nên hệ số đứng trước $\sin\omega_nt$ phải là $B_n/\omega_n$.

\textbf{Bước 4 — Ghép nghiệm và xét mode hằng.}


\[
u=\sum_n\left(A_n\cos\omega_nt+\frac{B_n}{\omega_n}\sin\omega_nt\right)X_n.
\]


Với Neumann hai đầu, thêm $A_0+B_0t$, trong đó $A_0=L^{-1}\int f$, $B_0=L^{-1}\int g$. Với tần số bị thay đổi bởi hạng $u$ hoặc cản $u_t$, giải lại ODE của từng mode; không dùng nguyên tần số $ak_n$.

\textbf{Bước 5 — Kiểm tra biên, dữ liệu và hội tụ.} Mỗi mode thỏa biên; tại $t=0$ hai chuỗi cho khai triển của $f,g$. Nếu dữ liệu không tương thích ở góc, giải theo nghĩa năng lượng hoặc từng miền và nêu đường mất tính trơn; không đạo hàm một chuỗi yếu tùy ý. Chu kỳ được suy từ tất cả mode có hệ số khác không, kể cả mode hằng; hai tần số có tỷ số vô tỷ ngăn một chu kỳ chung.

\textbf{Chi tiết tính tích phân hệ số.} Với $k=n\pi/L$, hai nguyên hàm thường dùng là


\[
\int x\sin kx\,dx=-x\cos kx/k+\sin kx/k^2,\qquad
\int x\cos kx\,dx=x\sin kx/k+\cos kx/k^2.
\]


Với đa thức $P$, ghi cả hai lần tích phân:


\[
\int_0^L P\sin kx\,dx
=\left[-P\cos kx/k\right]_0^L+
\left[P'\sin kx/k^2\right]_0^L-
\frac1{k^2}\int_0^L P''\sin kx\,dx.
\]


Chỉ bỏ hạng biên sau khi thế $0,L$ và điều kiện thực sự của $P$. Nếu dữ liệu chia đoạn, áp dụng công thức trên từng đoạn, rồi cộng các hạng tại điểm chia. Cuối cùng nhân nghịch đảo chuẩn của hàm riêng; nếu hệ số là của vận tốc, chia thêm tần số khi ghép nghiệm. Với mẫu số bằng $0$ trong tích–tổng, quay lại tích phân gốc để tính riêng.
\subsection*{Lời giải áp dụng cho bài}
\textbf{Câu (i): tìm một hàm riêng của toán tử có đạo hàm cấp một.} Đặt $f=e^{-x}(\pi\cos\pi x+\sin\pi x)$. Đạo hàm trực tiếp cho $f'=-(\pi^2+1)e^{-x}\sin\pi x$, nên biên $f'(0)=f'(1)=0$. Đạo hàm thêm cho $f''+2f'=-(\pi^2+1)f$. Vì vậy $u=fT$ dẫn tới $T''+(\pi^2+1)T=0$, $T(0)=1,T'(0)=0$:
\[u=f(x)\cos(\sqrt{\pi^2+1}\,t).\]
Các tính toán trên kiểm tra đồng thời PDE, biên Neumann và dữ liệu đầu.

\textbf{Câu (ii), bước 1: đổi tọa độ và khử đạo hàm cấp một.} $z=\log(1+x)$, $\ell=\log2$. Vì $\partial_x=e^{-z}\partial_z$, $\partial_x((1+x)^2u_x)=u_{zz}+u_z$. Đặt $u=e^{-z/2}v$; tính $u_{zz}+u_z=e^{-z/2}(v_{zz}-v/4)$. Nguồn $t/\sqrt{1+x}=te^{-z/2}$, nên $v_{tt}=v_{zz}-v/4+t$, biên Dirichlet ở $z=0,\ell$.

\textbf{Câu (ii), bước 2: chiếu và giải nguồn tuyến tính.} $v_0=e^{z/2}(e^z-1)(2-e^z)$, $v_{t,0}=0$. Với $k_n=n\pi/\ell$, $\omega_n=\sqrt{k_n^2+1/4}$,
\[A_n=\frac2\ell\int_0^\ell e^{z/2}(e^z-1)(2-e^z)\sin k_nz\,dz,\quad c_n=\frac{2(1-(-1)^n)}{n\pi}.\]
ODE $T_n''+\omega_n^2T_n=c_nt$ cho nghiệm riêng dữ liệu không $c_n(t/\omega_n^2-\sin\omega_nt/\omega_n^3)$. Vì vậy
\[u=e^{-z/2}\sum_{n\ge1}\left[A_n\cos\omega_nt+c_n\left(\frac t{\omega_n^2}-\frac{\sin\omega_nt}{\omega_n^3}\right)\right]\sin k_nz.\]
Tích phân của $A_n$ là tích phân sơ cấp: khai triển $v_0=-e^{5z/2}+3e^{3z/2}-2e^{z/2}$ rồi dùng
\[\int_0^\ell e^{pz}\sin k_nz\,dz=\frac{k_n(1-(-1)^ne^{p\ell})}{p^2+k_n^2}.\]
Thay $p=5/2,3/2,1/2$ cùng các hệ số $-1,3,-2$ cho hệ số tường minh.

\textbf{Câu (iii), bước 1: chiếu nguồn và vận tốc.} PDE $u_{tt}+2u_t=u_{xx}+\sin x$, biên không trên $(0,1)$, $f=0,g=1$. Với $k_n=n\pi$, $\omega_n=\sqrt{k_n^2-1}$, vận tốc có $c_n=2(1-(-1)^n)/(n\pi)$. Tích–tổng cho nguồn
\[F_n=2\int_0^1\sin x\sin k_nx\,dx=\frac{2(-1)^{n+1}k_n\sin1}{k_n^2-1}.\]
\textbf{Câu (iii), bước 2: giải ODE có cản.} $T_n''+2T_n'+k_n^2T_n=F_n$, $T_n(0)=0$, $T_n'(0)=c_n$. Nghiệm riêng tĩnh $F_n/k_n^2$ được bù hai dữ liệu đầu bởi phần thuần nhất:
\[u=\sum_{n\ge1}\left[c_ne^{-t}\frac{\sin\omega_nt}{\omega_n}+\frac{F_n}{k_n^2}\left\{1-e^{-t}\left(\cos\omega_nt+\frac{\sin\omega_nt}{\omega_n}\right)\right\}\right]\sin k_nx.\]
Thế $t=0$ kiểm tra hai dữ liệu; từng ODE kiểm tra PDE. $g=1$ không tương thích biên cố định ở góc, nên câu (iii) không có nghiệm cổ điển toàn bao đóng; nguồn cũng có vết không phù hợp nếu đòi cấp hai đến biên.

\textbf{Chi tiết đạo hàm ở (i).} Với $f=e^{-x}(\pi\cos\pi x+\sin\pi x)$, quy tắc tích cho
\[f'=e^{-x}[-\pi\cos\pi x-\sin\pi x-\pi^2\sin\pi x+\pi\cos\pi x]
=-(\pi^2+1)e^{-x}\sin\pi x.\]
Đạo hàm tiếp: $f''=(\pi^2+1)e^{-x}(\sin\pi x-\pi\cos\pi x)$. Vì $2f'=-2(\pi^2+1)e^{-x}\sin\pi x$, cộng cho $f''+2f'=-(\pi^2+1)e^{-x}(\sin\pi x+\pi\cos\pi x)=-(\pi^2+1)f$.

\textbf{Chi tiết đổi ẩn và tính hệ số ở (ii).} $u=e^{-z/2}v$ cho $u_z=e^{-z/2}(v_z-v/2)$, $u_{zz}=e^{-z/2}(v_{zz}-v_z+v/4)$. Cộng $u_z$ được $e^{-z/2}(v_{zz}-v/4)$. Vì $(e^z-1)(2-e^z)=-e^{2z}+3e^z-2$, dữ liệu mới là $-e^{5z/2}+3e^{3z/2}-2e^{z/2}$. Nguyên hàm
\[\int e^{pz}\sin k_nz\,dz=
\frac{e^{pz}(p\sin k_nz-k_n\cos k_nz)}{p^2+k_n^2}\]
được kiểm tra bằng đạo hàm tích. Tại $0,\ell$, $\sin k_nz=0$, $\cos k_n\ell=(-1)^n$, nên
\[A_n=\frac{2k_n}{\ell}\left[
-\frac{1-(-1)^n2^{5/2}}{k_n^2+25/4}
+3\frac{1-(-1)^n2^{3/2}}{k_n^2+9/4}
-2\frac{1-(-1)^n2^{1/2}}{k_n^2+1/4}\right].\]
Với nguồn $c_nt$, nghiệm riêng $c_nt/\omega_n^2$ có vận tốc đầu $c_n/\omega_n^2$; trừ $c_n\sin\omega_nt/\omega_n^3$ để bù, cho đáp ứng dữ liệu không.

\textbf{Chi tiết nguồn và hằng số của (iii).} Tích–tổng cho
\[F_n=\left[\frac{\sin((1-k_n)x)}{1-k_n}
-\frac{\sin((1+k_n)x)}{1+k_n}\right]_0^1
=(-1)^n\sin1\left(\frac1{1-k_n}-\frac1{1+k_n}\right).\]
Rút gọn cho $F_n=2(-1)^{n+1}k_n\sin1/(k_n^2-1)$. Đặt $H_n=F_n/k_n^2$; tổng nghiệm ODE là $T_n=H_n+e^{-t}(A\cos\omega_nt+B\sin\omega_nt)$. $T_n(0)=0$ cho $A=-H_n$; $T_n'(0)=c_n$ cho $-A+\omega_nB=c_n$, nên $B=(c_n-H_n)/\omega_n$. Thế lại và gom hai nhóm cho công thức ở bước 2.\end{document}